% 2de graad, Functies (wat een functie is, eerstegraadsfuncties, de functie c/x).
% Hiroshige: regen in rechte strepen met een vaste helling.
% Motief: de rand is een rechte met richtingscoëfficiënt 1/6; het trapje toont 6 opzij en 1 omhoog.
\def\omslagbestand{omslag-bui.jpg}
\def\omslagbron{Utagawa Hiroshige, \textit{Plotselinge bui boven de Shin-\=Ohashibrug en Atake}, 1857 (publiek domein)}
\def\omslagrand{10.2 + \x/6}
\def\omslagmotief{%
  \randlijn
  \pgfmathsetmacro{\ya}{10.2+13/6}%
  \draw[omslaglijn=1pt,dashed] (13,\ya) -- (19,\ya) -- (19,\ya+1);
  \fill[omslagstip] (13,\ya) circle (1.3mm);
  \fill[omslagstip] (19,\ya+1) circle (1.3mm);
  \node[text=white,font=\small,anchor=north] at (16,\ya-.1) {$+6$};
  \node[text=white,font=\small,anchor=west] at (19.1,\ya+.45) {$+1$};
  \node[text=fel,font=\large,anchor=north] at (16,\ya-.75) {$m=\frac16$};
}
